\documentclass[a4paper,11pt]{article}

\usepackage[margin=2.4cm]{geometry}
\usepackage[T1]{fontenc}
\usepackage{lmodern}
\usepackage{booktabs}
\usepackage{enumitem}
\usepackage{hyperref}
\usepackage{xcolor}
\usepackage{listings}

\hypersetup{colorlinks=true,linkcolor=black,urlcolor=blue}
\definecolor{codebg}{RGB}{247,247,247}
\lstset{
  basicstyle=\ttfamily\small,
  backgroundcolor=\color{codebg},
  frame=single,
  breaklines=true,
  columns=fullflexible,
  showstringspaces=false
}

\title{Physics Learning Agent\\Architecture and Technical Design}
\author{Physics Learning Agent contributors}
\date{\today}

\begin{document}
\maketitle

\begin{abstract}
Physics Learning Agent is a focused workspace for undergraduate physics. It combines bilingual conversational tutoring, original practice generation, structured course knowledge, account-scoped document retrieval, and multi-provider model access. This document describes the system boundaries, LangGraph workflow, memory model, retrieval pipeline, streaming protocol, persistence strategy, and the engineering layer that distinguishes the application from a direct language-model chat interface.
\end{abstract}

\tableofcontents

\section{Product Scope}

The application supports general physics, mathematical methods for physics, theoretical mechanics, electrodynamics, quantum mechanics, and thermodynamics and statistical physics. Its primary workflows are conversational learning, structured topic exploration, original practice generation, and question answering grounded in user-owned notes.

The interface is intentionally restrained. The main workspace uses a persistent conversation sidebar, an independently scrolling message region, and a fixed composer. Markdown, tables, code, and LaTeX are rendered for long-form technical answers.

\section{Technology Stack}

\begin{tabular}{p{0.28\textwidth}p{0.64\textwidth}}
\toprule
Layer & Implementation \\
\midrule
Web application & Next.js App Router, React, TypeScript \\
Styling & Tailwind CSS \\
Agent workflow & LangGraph with typed deterministic nodes \\
Document processing & LangChain text splitters and format-aware extraction \\
Retrieval & BM25, phrase, heading, metadata, and diversity scoring \\
Rendering & React Markdown, remark-math, rehype-katex, KaTeX \\
Testing & Vitest, Playwright, ESLint, production build checks \\
Persistence & Browser storage plus account-scoped local JSON and files \\
\bottomrule
\end{tabular}

\section{Request Lifecycle}

Each generation is bound to a tuple:

\begin{center}
\texttt{conversationId + assistantMessageId + requestId}
\end{center}

The client creates the identifiers before sending a request. The API validates request size and fields, applies a lightweight rate limit, resolves the signed-in user, and invokes the LangGraph workflow. The graph classifies the request, resolves course and topic context, updates structured memory, plans retrieval, retrieves relevant snippets when needed, and prepares the final request. The prompt builder then supplies task-specific instructions to the selected provider.

Provider-specific streaming formats are normalized by the server. The browser appends a chunk only when the generation tuple is still active. Conversation switches, new conversations, deletions, and explicit stops abort the request; late chunks are ignored.

\section{Agent Workflow}

The graph has six nodes:

\begin{enumerate}[leftmargin=2em]
  \item \textbf{Understand input}: intent, response language, practice style, and reference profile.
  \item \textbf{Resolve context}: course, topic, query type, knowledge mode, and bounded history.
  \item \textbf{Update memory}: learning goal, confusions, covered concepts, exercise direction, and preferences.
  \item \textbf{Plan retrieval}: decide whether personal or bundled knowledge is relevant.
  \item \textbf{Retrieve knowledge}: obtain a small set of account-scoped and optional sample snippets.
  \item \textbf{Prepare generation}: assemble a typed prompt-ready request.
\end{enumerate}

The workflow uses one primary model generation call. LangGraph provides explicit state management without introducing a high-latency multi-agent chain.

\section{Memory and Context}

Recent message history, conversation memory, and workspace-level learning preferences are separate. The prompt receives a bounded recent window and structured summaries rather than an unlimited transcript. The memory retains the active course and topic, current goal, recent confusions, covered concepts, practice direction, language, and preferred depth. Temporary provider keys are never included in synchronized workspace data.

\section{Personal Knowledge Retrieval}

User-owned documents are extracted into metadata-rich chunks. Text formats use LangChain splitters; PDF and office documents use structure-aware extraction where possible. Retrieval combines bilingual tokenization, BM25, exact phrase matches, heading relevance, course/topic metadata, and duplicate suppression. Selected snippets include source and locator metadata and are injected through citation-aware prompt instructions.

The local scorer accepts optional dense-vector scores. A production migration can therefore add PostgreSQL and \texttt{pgvector} while retaining the existing extraction, workflow, and prompt interfaces.

\section{Advantages over Direct Model Chat}

The application adds several deterministic layers that a direct model prompt does not provide:

\begin{itemize}[leftmargin=2em]
  \item course-aware terminology and answer structures for physics learning tasks;
  \item stable intent routing that avoids applying physics templates to unrelated questions;
  \item language-aware Chinese and English reference profiles;
  \item original practice generation with explicit difficulty, source style, and output controls;
  \item structured short-term and learning memory for follow-up questions;
  \item selective retrieval from the signed-in user's own materials;
  \item strict streaming isolation between conversations;
  \item continuation and retry behavior for interrupted long answers;
  \item consistent Markdown and LaTeX rendering and editable TeX export;
  \item provider abstraction for a server default or temporary Bring Your Own Key access.
\end{itemize}

These layers do not guarantee factual correctness. They improve context selection, task consistency, learning structure, and operational reliability. Users should still verify formal derivations against course materials.

\section{Persistence and Security Boundaries}

Anonymous state is stored in the browser. Signed-in state is synchronized to account-scoped files below \texttt{PLA\_DATA\_DIR}. Passwords use salted scrypt hashes and session cookies are HTTP-only, SameSite, and secure in production. Atomic file replacement and keyed in-process locks reduce corruption and lost updates in a single Node process.

Custom model URLs are checked on the server. Production endpoints require HTTPS and private, loopback, link-local, metadata, and reserved addresses are rejected. Trusted self-hosted private gateways require the explicit \texttt{PLA\_ALLOW\_PRIVATE\_MODEL\_ENDPOINTS=true} setting.

The local store and in-memory rate limiter are intended for development and small single-instance deployments. Multi-instance production requires transactional storage, shared sessions, shared rate limits, object storage, and background ingestion workers.

\section{Verification}

\begin{lstlisting}
npm run lint
npm run test:run
npm run build
npm run test:e2e
\end{lstlisting}

The suite covers classification, prompt construction, memory, stream guards, request limits, endpoint policy, persistence concurrency, retrieval, document parsing, conversation isolation, mobile layout, KaTeX, and knowledge-base flows.

\end{document}
